\documentclass[11pt]{article}
\usepackage[margin=1in]{geometry}
\usepackage[T1]{fontenc}
\usepackage{lmodern}
\usepackage{microtype}
\usepackage{graphicx}
\usepackage{amsmath,amssymb,amsthm,mathtools}
\usepackage{enumitem}
\usepackage[dvipsnames]{xcolor}
\usepackage[hidelinks]{hyperref}
\usepackage[nameinlink,capitalise]{cleveref}

\newtheorem{theorem}{Theorem}[section]
\newtheorem{lemma}[theorem]{Lemma}
\theoremstyle{remark}
\newtheorem{remark}[theorem]{Remark}

\newcommand{\K}{\mathcal K}
\newcommand{\F}{\mathcal F}
\newcommand{\C}{\mathbb C}
\newcommand{\czerosum}{\mathop{\bigoplus}\limits^{c_0}}

\title{A Non-Dual Positive Subcase of the 2-Local Triple-Homomorphism Problem\\
\large Compact C*-algebra domains}
\author{Solution packet for arXiv:1405.3836}
\date{22 July 2026}

\begin{document}
\maketitle

\begin{abstract}
Burgos, Fernandez-Polo, Garces, and Peralta ask whether every 2-local
triple homomorphism between C*-algebras is linear, and prove this when the
domain is a von Neumann algebra (more generally, a JBW*-triple).  We prove a
non-dual subcase: if $A=\czerosum_{i\in I}\K(H_i)$ is a compact C*-algebra
and $F$ is any JB*-triple, then every 2-local triple homomorphism $T:A\to F$
is linear and hence a triple homomorphism.  The mechanism is a row-column
factorization on rank-one operators followed by singular-value decomposition.
The general problems remain open.
\end{abstract}

\section{Source problem and status}

Every C*-algebra is a JB*-triple for
\[
 \{a,b,c\}=\tfrac12(ab^*c+cb^*a).
\]
A map $T:E\to F$, not assumed linear or continuous, is a \emph{2-local
triple homomorphism} if for every $x,y\in E$ there is a bounded linear triple
homomorphism $\Phi_{x,y}:E\to F$ satisfying
$\Phi_{x,y}(x)=T(x)$ and $\Phi_{x,y}(y)=T(y)$.

Problem 1.1 of \cite{BurgosEtAl} asks:
\begin{quote}
Is every 2-local triple homomorphism between C*-algebras automatically
linear?
\end{quote}
Problem 1.2 asks the same for arbitrary JB*-triples.  The source paper proves
linearity when the domain is a JBW*-triple.  The result below covers a class
of non-dual domains.  Here ``compact C*-algebra'' means a C*-algebra of the
form $\czerosum_{i\in I}\K(H_i)$.

\begin{center}
\includegraphics[width=0.94\linewidth]{figures/open_problem_crop.png}

\small Source PDF, page 2: Problem 1.1 and adjacent scope context.
\end{center}

\section{Proof intuition}

The source theorem cannot simply be applied to $\K(H)$: for
infinite-dimensional $H$, the compact operators are not a dual Banach space.
However, $\K(H)$ is densely assembled from rank-one operators.  If one vector
in a rank-one tensor is held fixed, the resulting row or column is a Hilbert
space and a JBW*-subtriple.  The source theorem therefore makes $T$ linear on
every row and every column.

After applying a functional $\varphi\in F^*$, these two families of linearity
statements say that
\[
 (\xi,\eta)\longmapsto \varphi\bigl(T(\theta_{\xi,\eta})\bigr)
\]
is sesquilinear.  It consequently defines a linear functional on finite-rank
operators.  A singular-value decomposition writes a finite-rank operator as
a linear combination of mutually orthogonal rank-one tripotents.  The finite
orthogonal-tripotent lemma in the source paper then identifies this functional
with $\varphi\circ T$ and, crucially, bounds it by the operator norm rather
than the trace norm.  It extends to $\K(H)$, proving that every scalarization
of $T$ is linear.  Finally, norm approximation by finitely many elementary
summands handles a $c_0$-sum.

\section{The elementary case}

We use the following consequences of 2-locality proved in
\cite[Section 2]{BurgosEtAl}.  For arbitrary JB*-triples $E,F$, every
2-local triple homomorphism $T:E\to F$ is complex homogeneous and
1-Lipschitz, and if $e_1,\ldots,e_n$ are mutually orthogonal tripotents, then
\begin{equation}\label{eq:orthogonal}
 T\left(\sum_{j=1}^n\lambda_j e_j\right)
   =\sum_{j=1}^n\lambda_jT(e_j)
 \qquad(\lambda_j\in\C).
\end{equation}
The same paper proves that a 2-local triple homomorphism whose domain is a
JBW*-triple is linear \cite[Theorem 3.8]{BurgosEtAl}.

For $\xi,\eta\in H$, let
\[
 \theta_{\xi,\eta}(z)=\langle z,\eta\rangle\xi
 \qquad(z\in H),
\]
using the convention that the inner product is linear in its first variable.

\begin{lemma}\label{lem:rows-columns}
Let $T:\K(H)\to F$ be a 2-local triple homomorphism.  For every unit vector
$\eta$, $T$ is linear on the column
\[
 C_\eta=\{\theta_{\xi,\eta}:\xi\in H\},
\]
and for every unit vector $\xi$, it is linear on the row
\[
 R_\xi=\{\theta_{\xi,\eta}:\eta\in H\}.
\]
\end{lemma}

\begin{proof}
Both $C_\eta$ and $R_\xi$ are closed JB*-subtriples of $\K(H)$.  For example,
\[
 \{\theta_{x,\eta},\theta_{y,\eta},\theta_{z,\eta}\}
 =\tfrac12\bigl(\langle z,y\rangle\theta_{x,\eta}
                 +\langle x,y\rangle\theta_{z,\eta}\bigr),
\]
so $C_\eta$ is the standard Hilbertian Cartan factor; the row calculation is
the same.  A Hilbertian Cartan factor is reflexive and is a JBW*-triple.

For two elements of $C_\eta$, restrict an implementing triple homomorphism
for $T$ on $\K(H)$ to $C_\eta$.  Thus $T|_{C_\eta}$ is itself a 2-local
triple homomorphism.  The JBW*-domain theorem
\cite[Theorem 3.8]{BurgosEtAl} makes it linear.  The same argument applies to
$R_\xi$.
\end{proof}

\begin{theorem}\label{thm:elementary}
Let $H$ be a complex Hilbert space and $F$ a JB*-triple.  Every 2-local triple
homomorphism $T:\K(H)\to F$ is linear and a triple homomorphism.
\end{theorem}

\begin{proof}
Fix $\varphi\in F^*$ and set
\[
 b_\varphi(\xi,\eta)
   =\varphi\bigl(T(\theta_{\xi,\eta})\bigr).
\]
By Lemma~\ref{lem:rows-columns}, $b_\varphi$ is linear in $\xi$ and
conjugate-linear in $\eta$.  Homogeneity deals with non-unit fixed vectors.
The Lipschitz estimate for $T$, together with $T(0)=0$, gives
\begin{equation}\label{eq:b-bound}
 |b_\varphi(\xi,\eta)|
 \leq \|\varphi\|\,\|\xi\|\,\|\eta\|.
\end{equation}
Hence $b_\varphi$ induces a well-defined algebraic linear functional
$L_\varphi$ on the finite-rank operators by
\[
 L_\varphi\left(\sum_{k=1}^m\theta_{\xi_k,\eta_k}\right)
   =\sum_{k=1}^m b_\varphi(\xi_k,\eta_k).
\]
Indeed, the canonical map from the algebraic tensor product
$H\otimes\overline H$ onto the finite-rank operators is an algebraic
isomorphism, and the sesquilinearity just proved is precisely the balancing
needed for this definition.

Let $x$ be finite-rank.  Choose a singular-value decomposition
\[
 x=\sum_{j=1}^n s_j\theta_{\xi_j,\eta_j},
\]
where $s_j\geq0$ and both $(\xi_j)$ and $(\eta_j)$ are orthonormal.  The
operators $e_j=\theta_{\xi_j,\eta_j}$ are mutually orthogonal rank-one
tripotents.  Consequently, \eqref{eq:orthogonal} yields
\[
 T(x)=\sum_{j=1}^n s_jT(e_j),
 \qquad
 L_\varphi(x)=\varphi(T(x)).
\]
It follows that
\begin{equation}\label{eq:operator-bound}
 |L_\varphi(x)|
   =|\varphi(T(x))|
   \leq\|\varphi\|\,\|T(x)\|
   \leq\|\varphi\|\,\|x\|.
\end{equation}
Thus $L_\varphi$ extends uniquely to a bounded linear functional on
$\K(H)$.  If $x_n$ are finite-rank and $x_n\to x$ in norm, continuity of
$T$ and $L_\varphi$ gives
\[
 \varphi(T(x))=\lim_n\varphi(T(x_n))
              =\lim_nL_\varphi(x_n)=L_\varphi(x).
\]
Therefore $\varphi\circ T$ is linear for every $\varphi\in F^*$.  Since
$F^*$ separates points, $T$ is additive; it was already complex homogeneous,
so it is linear.  Finally, a linear 2-local triple homomorphism preserves the
triple product by the polarization argument in
\cite[Section 2]{BurgosEtAl}.
\end{proof}

\section{Compact C*-algebras}

\begin{theorem}[compact-domain subcase]\label{thm:compact}
Let
\[
 A=\czerosum_{i\in I}\K(H_i)
\]
and let $F$ be a JB*-triple.  Every 2-local triple homomorphism
$T:A\to F$ is linear and a triple homomorphism.
\end{theorem}

\begin{proof}
Write $A_i=\K(H_i)$.  Restricting the implementing triple homomorphisms shows
that $T_i=T|_{A_i}$ is a 2-local triple homomorphism.  By
\cref{thm:elementary}, each $T_i$ is linear.

If $i\neq j$, $x\in A_i$, and $y\in A_j$, choose finite singular-value
truncations $x_n\to x$ and $y_n\to y$.  Each $x_n$ and $y_n$ is a finite
linear combination of rank-one tripotents, and all tripotents appearing in
the two different summands are mutually orthogonal.  Applying
\eqref{eq:orthogonal} and then passing to the norm limit gives
\begin{equation}\label{eq:cross-add}
 T(x+y)=T(x)+T(y).
\end{equation}
Induction gives the analogous identity for every finite set of distinct
summands.

For $a=(a_i)_{i\in I}\in A$, its finite coordinate truncations converge to
$a$ in norm.  Since $T$ is Lipschitz,
\[
 T(a)=\lim_{J}\sum_{i\in J}T_i(a_i),
\]
where $J$ ranges over finite subsets of $I$.  For $a,b\in A$, linearity of
every $T_i$ now gives
\[
 T(a+b)=\lim_J\sum_{i\in J}T_i(a_i+b_i)
       =T(a)+T(b).
\]
Together with complex homogeneity, this proves linearity.  The polarization
argument cited in \cref{thm:elementary} proves that $T$ is a triple
homomorphism.
\end{proof}

\section{Verification, limitations, and novelty check}

\paragraph{Internal checks.}
The proof uses only four inputs from the source paper: homogeneity,
1-Lipschitz continuity, linearity on finite combinations of mutually
orthogonal tripotents, and the JBW*-domain theorem.  The two points most
deserving independent review are the row/column restriction in
Lemma~\ref{lem:rows-columns} and the operator-norm estimate
\eqref{eq:operator-bound}.  Notice that the latter estimate is not inferred
from the sesquilinear-form bound \eqref{eq:b-bound}; it comes from 2-locality
and orthogonal singular vectors.  This is what forces extension to the
operator-norm completion $\K(H)$.

\paragraph{Scope.}
Infinite-dimensional $\K(H)$ is not a JBW*-triple, so
\cref{thm:elementary} is not a restatement of the source theorem.  The result
does not cover arbitrary C*-algebras, and it gives no bidual-extension
principle for a general 2-local map.  Problems 1.1 and 1.2 therefore remain
open outside the stated compact-domain class.

\paragraph{Bounded novelty check.}
On 22 July 2026 we searched the run's parsed arXiv corpus for the exact source
questions and for ``2-local triple homomorphism'', and performed bounded
web/arXiv searches combining that phrase with ``compact C*-algebra'',
``$K(H)$'', and ``real rank zero''.  The searches returned the source paper
and papers citing it, but no statement of \cref{thm:compact}.  The source
introduction mentions a compact-domain theorem for 2-local
\emph{*-homomorphisms}, not for triple homomorphisms.  This is evidence only;
the present deduction has not been peer reviewed and the novelty confidence
is moderate.

\paragraph{Recommendation.}
Promote as a likely-valid partial result, subject to expert review of the
three checkpoints listed in the accompanying README.  Do not count it as a
full answer to either source problem.

\begin{thebibliography}{9}
\bibitem{BurgosEtAl}
M. J. Burgos, F. J. Fernandez-Polo, J. J. Garces, and A. M. Peralta,
\emph{2-local triple homomorphisms on von Neumann algebras and
JBW*-triples}, J. Math. Anal. Appl. \textbf{426} (2015), 43--63;
\href{https://arxiv.org/abs/1405.3836}{arXiv:1405.3836},
doi:10.1016/j.jmaa.2014.12.058.
\end{thebibliography}

\end{document}
